\documentclass[10pt,a4paper]{article}
\usepackage[margin=2.6cm]{geometry}
\usepackage[T1]{fontenc}
\usepackage{lmodern,amsmath,amssymb,amsthm,microtype,booktabs,tabularx,listings,xcolor}
\usepackage[colorlinks,linkcolor=black,citecolor=black,urlcolor=black]{hyperref}
\lstset{basicstyle=\small\ttfamily,columns=fullflexible,keepspaces=true,
  breaklines=true,showstringspaces=false,frame=single,rulecolor=\color{black!20}}
\newtheorem{definition}{Definition}
\newtheorem{conjecture}{Conjecture}
\newcommand{\keep}{\ensuremath{\mathsf{keep}}}
\newcommand{\any}{\mathsf{any}}
\newcommand{\norm}{\mathsf{normal}}
\newcommand{\exc}{\mathsf{exceptional}}
\newcommand{\absent}{\mathsf{none}}
\newcommand{\some}{\ensuremath{\mathsf{some}}}
\newcommand{\code}[1]{\texttt{#1}}
\newcommand{\lean}[1]{\texttt{\small #1}}
\newcommand{\status}[1]{\textsf{\small[#1]}}
\newcommand{\ruleof}[3]{\frac{\begin{gathered}#2\end{gathered}}{#3}\;\textsc{#1}}
\newcommand{\respects}{\mathrel{\triangleleft}}
\newcommand{\supp}{\operatorname{supp}}
\newcommand{\proj}{\operatorname{project}}
\newcommand{\priv}{\operatorname{private}}
\newcommand{\slots}{\operatorname{slots}}
\newcommand{\du}{\operatorname{dedup}}
\newcommand{\incl}{\operatorname{covers}}
\newcommand{\ren}{\operatorname{rename}}
\newcommand{\at}{\operatorname{at}}
\newcommand{\argroots}{\operatorname{actuals}}
\newcommand{\spine}{\operatorname{spine}}
\newcommand{\termsame}{\mathrel{\#}}
\title{Inferred callback preservation\\[2pt]\large Effect requirements for heap-dependent refinements in TScert}
\author{TScert working paper}
\date{22 November 2026}

\begin{document}
\maketitle

\begin{abstract}\noindent
A function can depend on a heap refinement across an invocation of a
callback supplied by its caller. This paper describes the TScert
experiment that infers a preservation requirement from the function's
ordinary checking state and publishes it in the function's interface.
At an application, the checker compares that requirement with the
callback's inferred effect after substituting invocation arguments and
closure captures. The rules reuse the supports and separation queries
of \emph{Tolerant frames}. We give executable judgments and pseudocode
for the annotated source path, including its two body walks and direct
forwarding. Source tests exercise nonempty requirements in one file
and across module boundaries. The guarded Pretext walker produced an
empty requirement. A separate Lean control uses a length equality as
a premise. The source feature is experimental and its composition
with the verifier's soundness theorem remains open.
\end{abstract}

\section{The problem}\label{sec:problem}

\begin{lstlisting}
type Visitor = (i: number) => void;
export function walk(xs: number[], visit: InferPreserves<Visitor>): number {
  const c = { i: 0 };
  let sum = 0;
  /** @invariant safeInteger(c.i) && 0 <= c.i */
  while (c.i < xs.length) {
    visit(c.i);
    sum += xs[c.i]!;
    c.i++;
  }
  return sum;
}
\end{lstlisting}

The loop guard relates the cursor field to the length of the array.
The indexed read uses that relationship after the callback. The
callable type alone permits updates to the array. The local cursor is
allocated in the walker. In this example it is copied or accessed
through its own fields, and the callback receives a number. The
remaining requirement concerns the array supplied by the caller.

\code{InferPreserves<T>} has the TypeScript meaning \code{T}. TScert's
parser recognizes the declaration and records the position of the
marked parameter. The annotation requests inference. The interface
depends on the annotated function body, so functions sharing the
same TypeScript callback type may publish different requirements.

\emph{Tolerant frames}~\cite{tf} provides the ordinary rule for
retaining a fact across a step with a known effect. \emph{Read-only
frames}~\cite{ro} uses a component protocol and read-only reachability
to bound outside code. The judgment here uses a callback requirement
exported by the function that invokes it and checked where a caller
supplies it.

\section{Effects, supports, and requirements}\label{sec:metalanguage}

\begin{definition}[Footprint and support]
A footprint is a list of pairs $(r,\delta)$, where $r$ denotes a heap
region and $\delta$ is a bound on change. Write $\sqsubseteq$ for
\lean{Change.le}, the order
\[
\keep\sqsubseteq \mathsf{init}\sqsubseteq \mathsf{scalar}
\sqsubseteq \mathsf{extend}\sqsubseteq \mathsf{grow}\sqsubseteq \any.
\]
An effect footprint bounds writes. A support footprint bounds the
changes a fact tolerates. The partial function $\supp(a)$ used in
the source algorithm is \lean{IndexFact.support?} with the empty
read context. It is applied to the fact of an index atom.
\end{definition}

Regions include a field at a root and path, own fields, element
cells, array spines, class regions, restriction to a class, and
all locations. A root map $\pi$ assigns a root a destination
root and a path to the same value. $\proj_\pi$ is the exact partial
projection of a footprint. For a rooted region it prepends the
destination path, for a class restriction it projects the underlying
region, and for a rootless region it leaves it as written. It
normalizes empty paths for fields, own fields, and cells. A missing
root makes the projection undefined. Write $\ren_\sigma(\Phi)$ for
effect instantiation by \lean{Footprint.renameM}, whose map assigns
slots to slots. A missing slot widens the rooted region to all
locations, subject to any surrounding class restriction. These are
different operations because a demand identifies the locations a
fact depends on and an effect overapproximates writes.

\begin{definition}[Demand]
One inferred demand is a tuple
\[
 d=(p,S^{\norm},S^{\exc},A)
\]
where $p$ is a marked parameter's entry slot, $S^e$ is the support
required of that callback at exit $e$, and $A$ is a list of optional
entry-root paths for the receiver and ordinary arguments of the
callback invocation, in that order. An interface is a list of
demands. Its demands are conjunctive. An entry in the list names
an invocation through its parameter and actual-root list, so
different invocations of the same parameter can impose different
requirements.
\end{definition}

Given a separation query $K$, the existing stability check is
\[
 S \respects_K \Phi
 \quad\text{iff}\quad
 \forall(r,\delta)\in S.\ \forall(s,\gamma)\in\Phi.\quad
     \gamma\sqsubseteq\delta\ \lor\ K(r,s).
\]
Here $K(r,s)$ means the query reports the regions disjoint. This
is the executable \lean{respects} test, exposed by \lean{stable}.
The existing theorem \lean{respect\_sound} gives the support
condition from a sound separation query and a transition satisfying
the effect. Existing support lemmas then justify retaining the fact
under their premises. \status{Lean}

\section{Proof rules for the source checker}\label{sec:rules}

The new rules in this section specify experimental checking
judgments. They are implemented as decisions in Lean. The theorem
connecting the judgments to source-program safety is
stated separately in Section~\ref{sec:soundness}.

\subsection{Provenance and local discharge}

At each lowered label $\ell$, the table $O_\ell$ maps slots to
sets of possible origins. An origin is an input slot, an allocation
label, a field-read label, or an unknown origin. For the producer,
$\pi_\ell(x)=(j,\epsilon)$ precisely when
$O_\ell(x)=\{\mathsf{input}(j)\}$. A call is marked when the
instruction invokes a slot with a singleton input origin and that
input is marked. The private-allocation scan in
Section~\ref{sec:producer} produces a set $U$ of allocation labels.
Write $\priv_{O_\ell,U}(r)$ when $r$ selects an own field of a slot
with a singleton allocation origin in $U$, including the equivalent
empty-path field and class restrictions of either region. This
predicate captures the local discharge used by the source algorithm.

\subsection{Selecting and replaying facts}

Let $K_\ell$ be the knowledge materialized from the ordinary entry
state at a marked call. For an unconditional index atom $a$, write
$R_\ell(S)=\{(r,\delta)\in S\mid \neg\priv_{O_\ell,U}(r)\}$
for the residual support. Let $\mathsf{freeze}(P)$ replace every
change bound of $P$ by $\keep$.
\[
\ruleof{Candidate}
 {a\in K_\ell.\mathsf{indices}\qquad a.\mathsf{cond}=\absent
      \qquad \supp(a.\mathsf{fact})=S\ne[]\\
  \proj_{\pi_\ell}(R_\ell(S))=P}
 {K_\ell;O_\ell;U\vdash a\leadsto (S,\mathsf{freeze}(P))}
\]
Take all results of this rule. Their original supports form a
list $B$ and their projected requirements form $P_\ell$, with
duplicates removed from both lists. The demand is
$(p,P_\ell,[],A)$, where $A$ applies $\pi_\ell$ to each slot operand
of the invocation and records absence for other operands. An
atom with undefined support or undefined projection contributes
neither a retained support nor a requirement. This is a
conservative selection of available facts, with $\keep$ imposed
on every residual region. A separate backwards liveness or
minimality analysis is outside this judgment.

A walk checks the body using an ordinary step and the additional
supports $B$ at the marked invocation. On normal continuation,
\lean{framedKept} can retain an index atom from the step's input
knowledge if its computed support lies inside one of the supplied
supports. The test is inclusion of exact region--change entries.
The checker shifts the retained indices across the result slot and
inserts them into the post-step store.
\[
\ruleof{Retain}
 {a\in K_{\rm step}.\mathsf{indices}\qquad
  \supp(a.\mathsf{fact})=S\qquad \exists T\in B.\ S\subseteq T}
 {\mathsf{shift}(a)\in\mathsf{normalContinuation}(K_{\rm step},B)}
\]
This rule describes the additional facts inserted by
\lean{keptStep}. It leaves the observation and the ordinary
instruction and edge checks to the single pass. The source path
uses the ordinary exceptional transfer, which drops heap-dependent
facts at a handler.

In the replay, an inferred candidate must be covered by the
completed list $I$ from the proposing walk. Let $I[p,A]$ be the
concatenated normal footprints of all entries of $I$ with that
parameter and exactly that actual-root list.
\[
\ruleof{Replay}
 {d=(p,S^{\norm},[],A)\qquad
  \forall(r,\delta)\in S^{\norm}.      \exists(s,\gamma)\in I[p,A].\ s=r\land\gamma\sqsubseteq\delta}
 {I\vdash d\;\mathsf{covered}}
\]
This test compares regions syntactically. It does not invoke the
separation query. An empty requirement is covered vacuously.
The publication rule below uses the ordinary success test of each
walk. The second walk recomputes the candidates from its own states.

\[
\ruleof{Publish}
 {\mathsf{walk}(G,\absent)=I\qquad
  \mathsf{walk}(G,\some(I))\mathbin{\downarrow}\\
  \mathsf{checkUses}(G,I_{\rm deps},\mathsf{marked})=F}
 {I_{\rm deps}\vdash G\;\mathsf{publishes}\;\du(I\mathbin{+\!+}F)}
\]
The graph $G$ carries the initial checking state and the marked
inputs. A down arrow denotes successful completion.
Section~\ref{sec:producer} gives the walks and dependency order.

\subsection{Application and forwarding}

At a direct application of a declaration with a published
interface, let $\pi_C$ map the declaration's input positions
to caller roots and paths. The checker obtains the callback's own
effect $E$ from matching code entries (Section~\ref{sec:effects})
and constructs a slot substitution $\sigma$ for that code.
The substitution composes $A$ with $\pi_C$ for the receiver and
ordinary arguments, uses origins at closure creation for captures,
and leaves a slot unknown if it cannot be mapped to a caller
slot. Let $K_C$ be the caller's separation evidence.
\[
\ruleof{Use}
 {d=(p,S^{\norm},S^{\exc},A)\\
  \forall e\in\{\norm,\exc\}.\quad
     \proj_{\pi_C}(S^e)=T^e\ \land\ T^e\respects_{K_C}\ren_\sigma(E^e)}
 {\pi_C;K_C\vdash (d,E,\sigma)\;\mathsf{checks}}
\]
For an instantiated demand, failure to project the support or
to find a matching callback effect refuses the check. The source
driver skips this rule for an empty demand or an absent or
statically undefined callback actual. The ordinary typing check
still applies to the call.

When that actual is itself a marked input $q$ of a forwarding
declaration, let $\pi_F$ be the restriction of $\pi_C$ to paths
rooted at inputs of the forwarding declaration.
\[
\ruleof{Forward}
 {d=(p,S^{\norm},S^{\exc},A)\qquad \mathsf{actual}_C(p)=q
      \qquad q\;\mathsf{marked}\\
  \forall e\in\{\norm,\exc\}.\ \proj_{\pi_F}(S^e)=T^e}
 {\pi_F\vdash d\leadsto(q,T^{\norm},T^{\exc},\pi_F\circ A)}
\]
Composition on the list $A$ concatenates paths and leaves
an unknown or unmappable argument absent. Missing support roots
refuse forwarding. The forwarding declaration appends these demands
to the demands inferred from its own invocations.

\section{The producer algorithm}\label{sec:producer}

The producer starts from the lowered graph, point claims,
obligations, and marked parameter slots for an ordinary synchronous
member. Its prerequisites include the ordinary expression
obligations, closure code checks, and callback parameter types
whose expanded alternatives are callable or undefined and include
a callable alternative. The implementation declines a stated
regional bound. It sets up the single pass with the written
assertions and entry indices, an empty list of candidate
head facts, and the ordinary bound when the member has a
stated effect or a freshness promise.

\begin{lstlisting}
origins(graph, widths, start, inputCount):
    entry = all absent
    entry[start] = [{Input(j)} for j in 0 .. inputCount-1]
    transfer(label):
        if entry[label] is absent: return
        slots = entry[label]
        result = origin of a copied slot, Read(label) for a field read
                 of a slot, Allocated(label) for an allocation, Other otherwise
        out = prepend result to slots if the continuation binds it
        for target in normal successors(label):
            merge slotwise union of remap(out, target) into entry[target]
        if label has a handler:
            merge slotwise union of remap([Other] ++ slots, handler)
    return graphWorklist(transfer, entry)

privateAllocations(graph, entry):
    collect allocation labels at reached nodes
    for each reached node:
        copies expose no slots, except returns expose the copied slot
        field reads of a slot expose no slots
        field writes through a slot expose the operand slots of the value
        other operations expose every operand slot
        if operand enumeration fails, expose every current slot
        mark each allocation origin among those slots exposed
    return allocation labels with no exposure
\end{lstlisting}

The allocation classifier recognizes the lowered allocation
constructors and follows the body of a lowered let. Edge remapping
uses the target's slot map and destination width. The shared
\lean{graphWorklist} starts with all labels queued in reverse
order, snapshots a node and its outgoing labels, and requeues
changed labels that are absent from the queue. The provenance
analysis supplies the limit
$(n_{\rm inputs}+n_{\rm nodes}+2)(1+\max_\ell\mathsf{width}_\ell)$.
The worklist allows $1+n_{\rm nodes}\cdot\mathsf{limit}$ transfers.
Exhaustion returns failure.

\begin{lstlisting}
walk(initial, origins, privateLabels, marked, interface?):
    state = initial
    demands = []
    while the single-pass queue has a label:
        remove the first label
        supports = []
        if the node invokes a marked input by singleton provenance:
            entry = ordinary visitEntry(state, label)
            K = materialize all slots at entry
            supports, demand = Candidate selections at K
            if interface is present: require Replay coverage for demand
            append demand
        state = ordinary visit(state, label, framed = supports at this label)
    require queue empty and all reachable labels visited
    return dedup(demands)

produce(member, dependencyInterfaces):
    check code requests and prepare ordinary walk
    O = origins(member.graph)
    U = privateAllocations(member.graph, O)
    I = walk(initial, O, U, markedParameters, absent)
    walk(initial, O, U, markedParameters, some I)
    F = check uses of dependencyInterfaces, forwarding marked inputs
    return dedup(I ++ F)
\end{lstlisting}

The actual walk visits at most $n_{\rm nodes}+1$ queue entries.
Both walks start from the same initial state. The second
recomputes candidates and checks coverage against the first
walk's interface before calling \lean{SinglePass.visit}.
\lean{writeInterface} handles the final dependency check and
adds forwarded demands. It refuses direct self-reference and
groups with multiple members. The source driver resolves
dependencies before the interface is used, retains local results
in the prepared-file cache, and reads imported results from the
module interface. The publication entry contains the declaration
name and position, receiver flag, and demands, including the
per-exit supports and optional argument paths.

\section{Inferring and consuming closure effects}\label{sec:effects}

\lean{inferCodes} checks code requests in order, admitting the
previously accepted code entries. It runs the ordinary single pass
with grading enabled, checks the code entry conditions, and requires
an accepted code walk. For a synchronous closure, the inference
reads each recorded visit's step certificate. An asynchronous
closure receives the unrestricted footprint on both exits. A
summary component is either absent, meaning the analysis has no
contribution, or a footprint, including the empty footprint.

\begin{lstlisting}
inferEffect(codeRequest, visitRecords):
    links = regionTable of the code graph, starting with each input slot
            rooted at itself (receiver, ordinary arguments, captures)
    local = all absent
    for record in visitRecords:
        step = record.stepCertificate
        if step.scope is false: summary = (some top, some top)
        else:
            N, X = first recorded frame on each exit, defaulting to top
            normal = absent if result type is empty, else some rebase(N)
            exceptional = some rebase(X)
            if instruction is a typed own-field write:
                normal = its rebased field-write frame, if normal is present
                exceptional = some empty
            summary = (normal, exceptional)
        join summary into local[record.label], componentwise
    return compose(codeGraph, local)
\end{lstlisting}

Ordinary rebasing uses the call's position-to-slot map and the
closure entry links. The field-write case uses the slot of the
field receiver and the identity slot map. In either case it removes
\keep{} entries from the local effect before rebasing and supplies
no evidence of stable path agreement to \lean{rebaseAll}.
That operation performs the existing conservative widening
for regions without entry witnesses.

A join of two present summaries deduplicates the concatenation
of their entries. An absent operand contributes nothing at a join.
Sequencing has the same operation when both operands are present
and is absent otherwise. Distinct entries for the same region
and different changes are retained. The composition worklist
initializes the entry label to \some([]) and both exits to absent.
A node with a present incoming footprint and local summary
sequences the incoming footprint with each local exit footprint.
It joins the normal result into all normal successors, and into
the normal summary at a return. It joins the exceptional result
into the handler entry when one exists, or into the exceptional
summary otherwise. The worklist limit is two plus the sum of
lengths of all present local exit footprints.

\begin{lstlisting}
checkUses(member, publishedInterfaces, markedForwardingSlots):
    codes = inferCodes(member)
    O = origins(member.graph)
    bind each published declaration to a singleton input origin
        using source-map descriptors and the resolved import identity
    for node in member.graph:
        reject a use of a known declaration escaping direct application
        if node is an invocation of a known singleton input:
            resolve callee actual operands to caller roots or paths
            for each published demand:
                skip an empty demand or an absent or undefined callback
                if the callback actual is a marked forwarding input:
                    apply Forward and collect its result
                else:
                    identify the actual's graph closure allocation
                    join the summaries of all matching checked code rows
                    substitute invocation arguments and captures
                    apply Use with caller spine-separation evidence
    require a resolved direct call for every supplied interface
    return deduplicated forwarded demands
\end{lstlisting}

\lean{closure} identifies a graph closure by allocation origin
and matches code rows on arity, start, expressions, continuations,
and capture count. It joins their summaries before substitution,
reading an absent component as the empty footprint. Ordinary
actuals map to caller slots when their composed paths are empty.
Capture origins are taken at the allocation label. The collected
rest slot and unresolved actuals map to unknown. For support
projection, \lean{operandAt} can recover a field path from a
field-read origin if the backwards predecessor walk to the read
encounters only observation steps. Its \lean{unchanged} predicate
uses side-effect-free terms, reads through a slot, and allocations
with pure initializers. The caller constructs spine-separation
facts for distinct root identities when at least one is an
allocation identity. An input identity and an allocation identity
at label $\ell$ are encoded by $i$ and $n_{\rm inputs}+\ell$
respectively. The concrete Boolean check is \lean{Demand.check}.

For a consumer unit, the source driver checks the uses of the
completed interfaces and then runs the ordinary group check.
Success on this path is reported as an experimental result.
The ordinary verified judgment remains governed by the existing
soundness theorem.

\section{Examples and evidence}\label{sec:examples}

\begin{lstlisting}
export function sumAndTrace(xs: number[]): number {
  const visited: number[] = [];
  return walk(xs, i => { visited.push(i); });
}
export function sumAndAlias(xs: number[]): number {
  const visited = xs;
  return walk(xs, i => { visited.push(i); });
}
\end{lstlisting}

In \code{walk}, the cursor field can be discharged locally. The
observed nonempty demand includes \keep{} on the input array's
spine. The native call in each closure contributes an extension
of the array spine. In \code{sumAndTrace}, the capture maps to the
fresh allocation; separation satisfies \textsc{Use}. In
\code{sumAndAlias}, the capture maps to the same input as the
demand. Extension exceeds \keep{} and separation fails.
The source regression also exercises the corresponding
ordinary callback argument, and forwarding through another
marked declaration. It checks both layouts of the test program,
with the declarations in one file and with imports across a
module boundary. \status{source regression}

The guarded Pretext walker \code{walkPreparedLinesRaw} takes the
annotated parameter below. Its caller \code{findWholeLine} passes
\code{flow.data} and the callback.
\begin{lstlisting}
onLine?: InferPreserves<InternalLineVisitor>
\end{lstlisting}
The source run
published an empty requirement and the experimental client check
passed under that interface. This observation measures interface
production and use for that pair. A Lean control models the
heap-dependent refinement
\[
 \code{prepared.widths.length}
   = \code{prepared.segmentFlags.length}.
\]
Given the equality as a premise, the control checks its frame
condition against writes to the cache fields \code{wholeWidth},
\code{wholeStart}, and \code{wholeEnd}. That is a lowered-rule
control with supplied facts and effects. The source experiment
has yet to demonstrate a nonempty inferred preservation requirement
on a Pretext use. \status{source run and Lean control}

\section{Soundness and scope}\label{sec:soundness}

\begin{conjecture}[Soundness target]
For an admissible entry state of a producer, assume each marked
callback invocation meets every applicable demand of its
published interface at the relevant exit. If the producer
check succeeds, the retained refinements hold on the
continuation where they are used. If a caller supplies callbacks
whose code and effects satisfy the premises of the caller
check, the caller establishes those interface assumptions.
\end{conjecture}

The proof requires connecting allocation provenance and the
escape scan to local discharge, exact projection to location
identity, and closure summary inference and substitution to
actual heap transitions. It then needs a replay lemma connecting
the retained atoms to facts in the normal continuation, and an
interface composition lemma for forwarding and publication. The
existing respect and support theorems apply once those premises
have been established.

The normal source producer instantiates only the normal demand.
Closure inference and the consumer carry both exit components.
The Lean \lean{Site.infer} and \lean{Site.checkPublished} controls
also model both exits and accept a premise bounding possible
writes. They exercise the rule boundary with supplied facts and
evidence. They are separate from \lean{CallbackInferencePlan.atCall},
which selects index atoms from the single pass and performs the
local discharge above. Source inference for exceptional
continuations and simultaneous recursive interfaces remains open.

\section{Correspondence with Lean}\label{sec:correspondence}

The following definitions locate the judgments and pseudocode in
\code{lean/Verification/Source/Claims/}, except where indicated.

\begin{center}
\small
\begin{tabularx}{\linewidth}{@{}l>{\raggedright\arraybackslash}X@{}}
\toprule
Construction & Module and definitions \\
\midrule
Provenance and local discharge &
  \lean{CallbackOrigins.infer}, \lean{privateAllocations}, \lean{entryRoot};
  \lean{CallbackInferencePlan.ownPrivate} \\
\textsc{Candidate} and two walks &
  \lean{CallbackInferencePlan.atCall}, \lean{covers}, \lean{visit},
  \lean{walk}, \lean{infer} \\
\textsc{Retain} &
  \code{Programs/SinglePass.lean}:
  \lean{supportedIn}, \lean{framedKept}, \lean{keptStep} \\
Closure effects &
  \lean{CallbackEffectInference.inferCodes}, \lean{infer},
  \lean{step}, \lean{rebaseStep}, \lean{compose} \\
\textsc{Use} and \textsc{Forward} &
  \lean{CallbackPreservation.project}, \lean{Demand.check},
  \lean{Demand.forward}; \lean{CallbackInterfaceCheck.check},
  \lean{closure}, \lean{operandAt}, \lean{separation}, \lean{writeInterface} \\
Publication and source verdict &
  \code{Source/RequestCli.lean}: \lean{produceModuleCallback},
  \lean{PreparedFile.callbackInterface}, \lean{PreparedFile.check} \\
Controls &
  \code{tests/callback-preservation.mts}; \code{lean/Verification/Tests/}
  \code{CallbackPreservationControls.lean} \\
\bottomrule
\end{tabularx}
\end{center}

\begin{thebibliography}{9}
\bibitem{tf} TScert working note.
\emph{Tolerant frames: a modal account of fact preservation without ownership}.
11 November 2026.
\bibitem{ro} TScert working note.
\emph{Read-only frames: facts of a component across calls of code outside it}.
18 November 2026.
\end{thebibliography}

\end{document}
